
\section{Le programme empirique H1--H8 et ce que les témoins en laissent}
\label{sec:program}

Cette section résume, sans les reproduire, les estimateurs et les tests de la
version arXiv~v1. Elle sert à une chose : expliquer pourquoi, de toute la
théorie des sections~\ref{sec:bridge} à~\ref{sec:dynamics}, seul le canal de
dépendance passe dans le papier principal.

\subsection{Les huit hypothèses}

\begin{description}
\item[H1 (signature de régime).] Sous stress, $H^{dep}$ baisse, $\DD$ monte, et
donc $\JJ$ monte.
\item[H2 (compensation partielle).] $\Delta H^{dep}_t=-\beta\,\Delta H^{vol}_t+u_t$
avec $\beta>0$, et $\beta_{\text{stress}}>\beta_{\text{calme}}$.
\item[H3 (asymétrie des sauts).] Les variations de $\JJ$ sont unilatérales,
au-delà de ce que produit le même estimateur sur des données sans structure
temporelle.
\item[H4 (relaxation).] $\JJ$ revient à la moyenne avec une demi-vie plus longue
que la mémoire de l'estimateur.
\item[H5 (signature d'asymétrie).] Les sauts positifs de $\JJ$ coïncident avec
une asymétrie négative du facteur de marché (postulat~P\ref{p:sign}).
\item[H6 (couplage des canaux).] Dépendance et forme sont faiblement couplées en
marché calme, fortement sous stress.
\item[H7 (pouvoir prédictif).] L'entropie d'aujourd'hui améliore, hors
échantillon, la prévision du risque à 21~jours par rapport à la volatilité.
\item[H8 (composition du budget de crise).] Regroupés en épisodes, les sauts ont
plus souvent le canal de dépendance pour contributeur dominant qu'un
rééchantillonnage sans mémoire.
\end{description}

\paragraph{Estimation, en une phrase.} Fenêtres glissantes de 252~jours,
covariance rétrécie à la \citet{ledoit2004}, déficit marginal par
l'approximation de \citet{hyvarinen1998} plutôt que par le développement
d'Edgeworth~\eqref{eq:edgeworth} (inutilisable sur des rendements quotidiens),
régime de stress défini par un seuil de volatilité calculé sur le passé
seulement, sauts détectés par un filtre à quatre sigmas sur $\Delta\JJ$.

\subsection{Pourquoi un témoin calibré est indispensable}

Le point théorique qui gouverne toute la lecture des résultats est le suivant.
$\JJ$ est une fonctionnelle \emph{convexe et positive} des moments d'ordre deux à
quatre estimés sur une fenêtre. Quand une observation à queue épaisse entre dans
la fenêtre, les moments estimés sautent, et la fonctionnelle convexe saute
\emph{vers le haut} quel que soit le signe de l'observation. Quand
l'observation vieillit (pondération exponentielle) ou sort de la fenêtre, la
fonctionnelle redescend graduellement. L'appareil de mesure fabrique donc, à
lui seul, la dent de scie que prédisent la proposition~\ref{prop:conditioning}
et le postulat~P\ref{p:relax}. Une simulation gaussienne i.i.d., sans aucun
événement, produit déjà une demi-vie apparente de 35~jours et 37~sauts positifs
contre aucun négatif au filtre à quatre sigmas.

D'où deux témoins, construits en rééchantillonnant la matrice de rendements
observée et en faisant tourner la chaîne identique deux cents fois : le témoin
\emph{i.i.d.} (lignes tirées indépendamment, queues marginales et dépendance
transversale conservées, toute chronologie détruite) et le témoin \emph{par
blocs} de 21~jours (qui conserve en outre le regroupement de volatilité). Une
statistique qui franchit les deux est une preuve ; une statistique qui en
franchit un seul n'est pas identifiée.

\subsection{Le verdict}

\begin{table}[htbp]
\centering
\small
\begin{tabular}{@{}llrrrl@{}}
\toprule
& Statistique & Données & i.i.d. & blocs-21 & Verdict \\
\midrule
H1 & écart de régime $H^{dep}$ & $-0{,}99$ & $-0{,}21$ & $-0{,}51$ & \textbf{franchit les deux} \\
H1 & écart de régime $\JJ$ & $1{,}93$ & $2{,}14$ & $1{,}97$ & échoue aux deux \\
H1 & écart de régime $\DD$ & $0{,}27$ & $1{,}25$ & $0{,}65$ & à contre-sens \\
H2 & $\beta$ (stress) & $0{,}20$ & $0{,}05$ & $0{,}11$ & échoue aux deux \\
H3 & sauts positifs & $158$ & $148{,}5$ & $133{,}2$ & blocs seulement \\
H3 & asymétrie de $\Delta\JJ$ & $12{,}26$ & $14{,}80$ & $12{,}97$ & à contre-sens \\
H4 & demi-vie (jours) & $120{,}8$ & $92{,}1$ & $162{,}8$ & entre les deux \\
H4 & étendue de $\JJ$ & $18{,}95$ & $10{,}05$ & $13{,}77$ & \textbf{franchit les deux} \\
H5 & $\mathrm{corr}(\Delta\JJ,\Delta S_{mkt})$ & $-0{,}24$ & $-0{,}04$ & $-0{,}07$ & échoue aux deux \\
H6 & écart de couplage & $0{,}59$ & $0{,}77$ & $0{,}67$ & à contre-sens \\
H7 & gain hors échantillon & \multicolumn{3}{c}{aucun sur deux panels} & rejetée \\
H8 & part d'épisodes à dominante dépendance & $46\,\%$ & $19$--$33\,\%$ & n.i. & i.i.d.\ seulement \\
\bottomrule
\end{tabular}
\caption{Verdicts face aux témoins (moyennes de deux cents réplications),
panel principal de vingt titres du S\&P~500, 1990--2022. Source :
\texttt{paper/tables/tab\_placebo\_sp500.tex} et sections H7--H8 de la v1.
n.i.\ : non identifié, à toute longueur de bloc testée entre 5 et 126~jours.}
\label{tab:program}
\end{table}

Deux statistiques franchissent les deux témoins sur le panel principal, et une
seule se reproduit sur les 48~portefeuilles sectoriels de Kenneth French :
l'écart de régime du canal de dépendance ($-3{,}65$~nats contre $-0{,}56$ et
$-2{,}46$). La phase~A l'a ensuite durci : il tient en version dévolatilisée
(Forbes--Rigobon) et en version en rangs, cette dernière étant la plus forte
face au témoin par blocs ($p=0{,}03$ sur le S\&P~500, $0{,}01$ sur les secteurs),
alors que le ratio d'absorption de \citet{kritzman2011} échoue au témoin par
blocs sur les deux panels.

\subsection{Ce que le verdict dit de la théorie}

Il faut séparer ce que le verdict touche et ce qu'il ne touche pas.

\paragraph{Ce qui n'est pas touché.} Les énoncés des
sections~\ref{sec:bridge} et~\ref{sec:channels} sont des identités et des
théorèmes. La proposition~\ref{prop:conditioning} aussi. Aucun résultat empirique
ne peut les réfuter ; il peut seulement montrer que les estimateurs disponibles
ne permettent pas de les voir à l'\oe uvre.

\paragraph{Ce qui est touché.} Les deux postulats. P\ref{p:relax} n'est soutenu
que par l'amplitude (l'étendue de $\JJ$), pas par un taux de relaxation : la
demi-vie se situe entre les deux témoins et son niveau dépend du témoin
choisi. P\ref{p:sign} n'est pas établi (H5). Le couplage des canaux (H6), que la
section~\ref{sec:dynamics} dérivait du caractère commun de la nouvelle, est
reproduit et même dépassé par des données rééchantillonnées.

\paragraph{Pourquoi le canal de dépendance survit.} Une explication
vraisemblable, cohérente avec la phase~A, tient à la forme de l'estimateur.
$\tfrac12\log\det R=\tfrac12\sum_k\log\lambda_k$ pèse tout le spectre de la
matrice de corrélation, y compris les petites valeurs propres, qui s'effondrent
quand la coupe transversale se couple sur quelques modes collectifs. Un seul
gros rendement isolé déplace peu ces petites valeurs propres. À l'inverse, le
canal de forme $\DD$ réagit fortement à une seule observation extrême, ce qui
est exactement la signature que reproduit le témoin. Les mesures de
concentration (ratio d'absorption, part de la première valeur propre) ne
regardent que le haut du spectre, où le témoin par blocs les rattrape.

\paragraph{Conséquence pour le papier principal.} L'indicateur retient le seul
canal identifié, en abandonnant $\DD$ et donc $\JJ$. La théorie des
sections~\ref{sec:channels} et~\ref{sec:dynamics} reste la justification de
\emph{pourquoi} ce canal mesure une information partagée (la
multi-information gaussienne, équation~\eqref{eq:ceiling} et annexe~\ref{app:gauss}),
mais la dynamique postulée n'est pas utilisée.
